\section{Introduction}
\label{sec:introduction}

This document collects design notes for VSC HVDC transmission. Each topic is
worked out here first, with its derivations, assumptions and sources; the
corresponding slides of the HVDC design presentation are then derived from
it.

Section~\ref{sec:creepage} covers the creepage distance of insulators for DC
overhead lines: why pollution, not overvoltage, sets it, how the site
pollution is measured, and how the measured pollution, the insulator type and
the altitude lead to the creepage of a 400\,kV DC line.

Section~\ref{sec:dc-fault} describes the DC fault ride-through of a
half-bridge MMC symmetrical monopole in two parts: first the principles ---
the two fault types, the recovery sequence, its switchgear, fault zones and
selectivity, the ratings it affects and the design alternatives --- and then
their application to one installation, the Pugalur--Thrissur link, with its
specified times and recorded recoveries.

Section~\ref{sec:interference} treats the interference a converter station
causes above the harmonic range: the power line carrier band and the radio
interference band, their regulatory position in the United States, the form
in which a specification states a verifiable requirement, and indicative
limits adapted to the elevation, climate and line density of the western
United States.

The numerical values given are indicative ranges from the standards and
published guidance. They are not agreed AmWatt design constants; a project
value follows from the measured site pollution and the supplier's test data.
